\section{Held-Out Test Evaluation}
\label{sec:held_out_performance}

After the validation decision is fixed, all three finalists are refitted on the combined training and validation samples and evaluated once on the eight-quarter test period. The purpose of retaining all three candidates in this comparison is to assess whether the validation conclusions generalize, not to reopen model selection.

\begin{table}[!htbp]
\centering
\footnotesize
\caption{Held-out test performance of the no-information benchmark and frozen finalist models.}
\label{tab:finalist-test-performance}
\begin{tabular}{lrrrr}
\toprule
Model & $\mathrm{MAE}$ & $\mathrm{RMSE}$ & $R^2$ & $\overline{\rho}$ \\
\midrule
No-information & 0.1302 & 0.1819 & -0.053 & --- \\
Ridge & 0.0937 & 0.1297 & 0.465 & 0.727 \\
XGBoost & 0.0912 & 0.1258 & 0.497 & 0.738 \\
Temporal GraphSAGE & 0.0898 & 0.1268 & 0.488 & 0.740 \\
\bottomrule
\end{tabular}
\end{table}


\noindent
Table~\ref{tab:finalist-test-performance} shows that all three models improve substantially on the no-information benchmark. Temporal GraphSAGE has the lowest test MAE, XGBoost has the lowest RMSE and highest $R^2$, and their rank correlations are nearly identical. Ridge is weaker but remains close. The held-out evidence therefore confirms the validation-stage conclusion: nonlinear and temporal graph models provide modest gains, but no candidate dominates consistently across metrics.

The quarterly diagnostics in Appendix~\ref{app:quarterly_test_performance} show that the models also move together through time. Their absolute errors vary across quarters, whereas their within-quarter rank correlations remain strong across the test period. The selected XGBoost specification is not replaced on the basis of these test outcomes.

Figure~\ref{fig:finalist-test-deciles} in Appendix~\ref{app:quarterly_test_performance} adds an outcome-severity perspective. The three error profiles remain close through the first nine realized-drawdown deciles and rise sharply in the tenth. The main remaining limitation is therefore not the choice among these three models, but the difficulty of predicting the magnitude of the most extreme downside outcomes.
